\documentclass[11pt,a4paper]{article}
\usepackage[utf8]{inputenc}
\usepackage[T1]{fontenc}
\usepackage[english]{babel}
\usepackage{amsmath,amssymb}
\usepackage[round]{natbib}
\usepackage{booktabs}
\usepackage{graphicx}
\usepackage[margin=2.5cm]{geometry}
\usepackage{hyperref}
\usepackage{xcolor}
\hypersetup{colorlinks=true,linkcolor=blue!50!black,citecolor=blue!50!black,urlcolor=blue!50!black}
\usepackage{caption}
\captionsetup{font=small,labelfont=bf}

\title{\bfseries Perfect Abstention Is Zero Knowledge:\\
Two Stable Fixed Points of a Self-Referential Abstention Objective\\
in a Small Language Model with Versioned Memory}

\author{Maximiliano Speranza\\
\small Independent Researcher, Buenos Aires, Argentina\\
\small \href{https://orcid.org/0009-0005-0413-8554}{ORCID: 0009-0005-0413-8554}\\
\small \texttt{maximiliano.speranza@gmail.com}}

\date{August 2026}

\begin{document}
\maketitle

\begin{abstract}
\noindent
A hallucination detector should predict its own model's errors, so the natural training target is
self-referential: label each example by whether the model \emph{would be wrong if it answered}. We
train that target in a $863{,}730$-parameter language model with a co-trained versioned archive,
holding architecture, data and budget fixed, and find that the objective has \textbf{two stable fixed
points}. Nine units split $5$--$4$ between them, and the split is governed by \textbf{how much
retrieval the model already had when abstention training began}: six units that started from scratch
split $2$--$4$, while three that inherited $12{,}000$ steps of retrieval pre-training (RECUP
$0.77$--$0.80$) reached the useful point $3$ of $3$.

The degenerate one is the finding. Four of nine units converge to a state that scores
$\texttt{nose} = 1.0000$ and $\texttt{invento} = 0.0000$---\emph{perfect} on both metrics the
abstention literature reports---while their global accuracy sits at $0.4065$, which is exactly the
trivial floor obtained by abstaining on everything. They are not detectors that fail. They are
\textbf{perfectly calibrated detectors of a generator that knows nothing}, and the calibration is
honest: the model does not recover, and it correctly reports that it does not. Their binary head has
collapsed to a constant with $\mathrm{AUC}\ 0.522$--$0.578$ on its own target, and that constant sits
where the prior predicts it: for the youngest unit, $\mathrm{logit}$ of the base error rate is
$1.513$ against a measured median of $1.507$.

The bifurcation is decided early and is sharply legible. At step $2500$, on the $512$-sample
evaluation batch the training loop already computes, the five surviving units had emitted
$2$, $4$, $4$, $8$ and $1$ answers; the four degenerate ones had emitted \textbf{zero}. The rule
``zero versus non-zero at step $2500$'' separates \textbf{40 of 40} runs with a binary head across
twelve families in this repository, and one unit survived on a \emph{single} sample out of $512$.
The early mute phase itself belongs to the self-referential target: on the paired control that
shares seed \emph{and} initialization, three units with a fixed, model-independent target read
$0.06$--$0.20$ against $0.98$--$0.99$ for their self-referential twins.

The degenerate point is \textbf{absorbing}, not merely slow. Measuring recovery inside one unit with
paired samples across $4000$ further steps gives $-0.0021$, or $0.4\sigma$: it does not move. We
also report a correction to our own first reading, in which a between-unit trend that survives a
noise control ($r = 0.986$) was nonetheless not extrapolable, because the curve is concave and a
linear fit hid a saturation.

The practical consequence is a metric claim, not an architectural one: \textbf{any success criterion
that reads abstention rate and fabrication rate without reading global accuracy is maximized by a
model that has learned nothing}, and a self-referential objective supplies a training dynamic that
reaches that maximum on its own, in a substantial fraction of runs, with no warning from either
metric.
\end{abstract}

\vspace{0.5em}
\noindent\textbf{Keywords:} abstention; hallucination detection; selective prediction; self-referential
objectives; training dynamics; degenerate solutions; small language models; persistent memory

\section{The metric a silent model wins}
\label{sec:intro}

The literature on abstention reports two numbers. One is the rate at which a model correctly declines
to answer questions it cannot answer; the other is the rate at which it fabricates. Both are
desirable, both are bounded, and a system that improves on both is normally taken to have improved.

A model that abstains on every question scores perfectly on both.

This is not a subtle observation, and it is not new as an abstract point: the selective-prediction
literature has always paired risk with coverage precisely because risk alone is trivially minimized
\citep{elyaniv2010selective, geifman2017selective}. What we report is that the degenerate corner is
not only reachable in principle but is an \textbf{attractor of a natural training objective}, that a
substantial fraction of runs land in it, and that once landed the state is stable under further
training.

The objective in question is the one a hallucination detector actually wants. A detector is useful if
it predicts \emph{this model's} errors, so the target should not be ``does the answer exist?'' but
``will this model be wrong if it answers?''. That target is available for free during training: take
the model's own \texttt{argmax}, compare it against the label, and use the comparison as a binary
target with the gradient stopped. It is the natural thing to do and, in this bench, it is what
produced the units this paper is about.

Its defect is structural rather than incidental. The label is a function of the model's own state, so
it moves as the model learns, and the map has more than one consistent solution. At initialization
the model is wrong on nearly everything, so the target is the constant~$1$ and the head learns to
predict ``I will be wrong''---which is \emph{true}. Escaping requires the generator to start being
right before the head has finished settling on the constant. It is a race between two learning
timescales, and \S\ref{sec:mudas} shows that both the seed and the amount of retrieval the model
already has when the race starts decide which one wins.

Both outcomes are self-consistent, and that is what makes the degenerate one stable:

\begin{center}
\begin{tabular}{ll}
\toprule
fixed point & what the model asserts, and whether it is true \\
\midrule
useful & ``I retrieve well, and I know when I will fail'' --- true \\
\textbf{degenerate} & ``I retrieve almost nothing, and I know that I retrieve almost nothing'' --- \textbf{also true} \\
\bottomrule
\end{tabular}
\end{center}

The degenerate model is not miscalibrated. It is honest, useless, and invisible to the two metrics
the field reports.

\paragraph{Contributions.}
\begin{enumerate}
\item We characterize the degenerate fixed point of a self-referential abstention objective and show
it is reached by $4$ of the $6$ units trained from scratch, and by none of the $3$ that inherited
retrieval pre-training (\S\ref{sec:mudas}).
\item We show the degenerate units are \emph{not} broken detectors: their heads have collapsed to the
prior, and we verify this quantitatively against the closed-form constant, not just qualitatively
(\S\ref{sec:prior}).
\item We show the bifurcation is decided by step $2500$ and is separable by a zero/non-zero rule on
data the training loop already produces, validated on $40$ runs (\S\ref{sec:bifurcacion}).
\item We show the degenerate point is absorbing rather than slow, using the first within-unit
measurement of retrieval in this line of work (\S\ref{sec:absorbente}).
\item We correct a first reading of our own data in which a real between-unit trend was extrapolated
with the wrong functional form (\S\ref{sec:correccion}).
\end{enumerate}

\section{Related work}
\label{sec:related}

\paragraph{Selective prediction.} The risk--coverage framing
\citep{elyaniv2010selective, geifman2017selective} exists because abstention trades against
usefulness, and SelectiveNet \citep{geifman2019selectivenet} makes the trade explicit by training a
selection head against a coverage constraint. Our degenerate units are the coverage constraint's
reason for existing, observed as a training outcome rather than as an argument.

\paragraph{Learning to say ``I don't know''.} Adding an \texttt{[IDK]} token to the vocabulary
\citep{cohen2024idk} and adding a separate binary head are the two dominant interfaces; we compared
four of them under matched budget in prior work \citep{speranza2026abstention} and found the head
superior in this regime. That comparison used a \emph{fixed} target---``does the answer exist?''---
and none of its units entered the state described here. The present paper is what happens when the
target is made self-referential instead.

\paragraph{Self-referential and bootstrapped targets.} Training a predictor against labels derived
from the current model is standard in self-training and pseudo-labelling, and the confirmation-bias
failure mode is well documented \citep{arazo2020pseudo}. Our contribution is not that
self-referential targets can collapse, which is expected, but the specific shape of the collapse in
the abstention setting: the collapsed state is \emph{correct}, scores perfectly on the standard
metrics, and is therefore not detectable as a failure by them.

\paragraph{Degenerate solutions and collapse.} Representation-learning methods that predict their own
outputs collapse to constants without an explicit remedy \citep{grill2020byol, chen2021simsiam}. The
mechanism here is the same family---a target that the model can satisfy by moving the target---but
the collapsed solution is not a constant representation, it is a constant \emph{decision}, and its
constancy is exactly what makes it score well.

\section{The bench}
\label{sec:bench}

The bench is the one described in \citet{speranza2026abstention}; we restate what is needed to read
the results and refer there for the rest.

\paragraph{Model.} A language model trained \emph{from scratch}, no pretrained component:
$863{,}730$ parameters, $3.5$~MB, closed vocabulary of $242$ human-readable tokens. A delta rule
recurrence plus a \textbf{co-trained persistent archive} with a learned order stamp in the key, read
by softmax and injected at block~0. The archive persists across sessions while the recurrent state is
reset between them.

\paragraph{Task.} Facts are stated in one session, corrected in another, and queried in a third. The
answer is a \textbf{single token}, so the metric is exact and deterministic: no LLM judge, no
interpretive parser. $40\%$ of questions have no answer ($p_{\mathrm{nose}} = 0.4$); the empirical
fraction on the evaluation sample is $0.4065$.

\paragraph{Interface.} All units in this paper use the separate binary head, which the prior
comparison selected: a scalar $a$ projected from the same final state that feeds the vocabulary
softmax, $+129$ parameters. The model abstains when $a > 0$. Reporting the collapse on the
\emph{winning} interface is deliberate---it is not an artifact of a weak abstention mechanism.

\paragraph{Metrics.} \texttt{vigente} is accuracy on the current value \emph{as emitted}, i.e.\ after
the head's decision; \texttt{nose} is the abstention rate on questions without an answer;
\texttt{falsa\_abst} is the abstention rate on questions that do have one; \texttt{invento} is the
rate of answers naming content absent from the archive. We add two instruments:

\begin{itemize}
\item \textbf{Global accuracy} $= (\text{correct answers} + \text{correct abstentions}) / n$. Its
trivial floor on this bench is $\mathbf{0.4065}$, obtained by abstaining on everything.
\item \textbf{RECUP}, retrieval accuracy \emph{ignoring the head}: the \texttt{argmax} over values on
questions that do have an answer, with the abstention decision not applied. \texttt{vigente} is a
composite of retrieval and decision and is $0.0000$ whenever $a > 0$ everywhere, regardless of what
the model retrieved; RECUP separates the two.
\end{itemize}

\paragraph{Measurement protocol.} RECUP is measured with $n = 8000$ and a fixed data seed
($54321$) across every checkpoint, so comparisons are \textbf{paired}: the questions are identical
and only the model differs. This matters. Measured across six data seeds at $n=2000$, RECUP has
$\sigma = 0.0135$, so an unpaired difference carries $\sigma \approx 0.0191$---larger than the
effects of interest. We report $p_{\mathrm{flip}}$, the fraction of questions on which two
checkpoints disagree, and derive the paired standard deviation from it rather than assuming a value.

\paragraph{Protocol.} Every campaign has a pre-registration frozen with its SHA before implementation,
a deviations file, and per-unit reporting. \textbf{Results are never reported as a mean alone.}
Bimodality across seeds is a measured property of this bench, and \S\ref{sec:mudas} is a direct
consequence of it.

\paragraph{Use of a language model.} All experiment code, training campaigns and report generation in
this work were carried out with the assistance of a large language model. The study design, the
pre-registrations and their frozen hashes, the choice of controls and the verification of every
reported number are the author's own. The model is not an author and bears no accountability for the
work.

\section{The objective}
\label{sec:objetivo}

Write $\ell$ for the value logits, $y$ for the label, $a$ for the head's scalar. The self-referential
target is
\[
b \;=\; \mathbb{1}\!\left[\operatorname{sg}\!\big(\arg\max\nolimits_{v \neq \texttt{NOSE}} \ell_v\big) \neq y\right],
\qquad
\mathcal{L} \;=\; \mathrm{BCE}(a, b) \;+\; \mathrm{CE}_{\text{value}},
\]
with $\operatorname{sg}$ the stop-gradient. On questions with no answer, $b = 1$ by construction.
The value cross-entropy is computed only where an answer exists and normalized by that fraction, so
that the loss scale does not depend on $p_{\mathrm{nose}}$.

Two properties follow directly and are worth stating before any measurement, because they are what
the results confirm.

\paragraph{The value gradient is never gated by the head.} $\mathrm{CE}_{\text{value}}$ does not
depend on $a$. Whatever happens to the head, retrieval keeps receiving gradient. The collapse
reported below is therefore \emph{not} an instance of ``the model stops learning because it stopped
answering''---a hypothesis we held and disproved (\S\ref{sec:mudas}).

\paragraph{If the head ignores its input, its optimum is the prior.} A constant $a$ minimizing
$\mathrm{BCE}$ against $b$ is $\mathrm{logit}\,\mathbb{E}[b]$, and $\mathbb{E}[b]$ is computable in
closed form from RECUP:
\begin{equation}
\mathbb{E}[b] \;=\; p_{\mathrm{nose}} + (1 - p_{\mathrm{nose}})\,(1 - \mathrm{RECUP}),
\qquad p_{\mathrm{nose}} = 0.4065 .
\label{eq:prior}
\end{equation}
This is a prediction with no free parameters, and \S\ref{sec:prior} tests it.

\section{Results}

Nine units, seeds $0$--$8$, identical flags. Five reach a useful state; four do not. We report per
unit throughout.

\paragraph{A confound in the starting point, found late and reported in full.} The nine units are
\emph{not} homogeneous, and we discovered this after the analysis below was written. The training
harness seeds the abstention phase from a base checkpoint when one exists for that seed. Only three
such bases existed (\texttt{n3\_s0}, \texttt{n3\_s1}, \texttt{n3\_s2}, each carrying $12{,}000$
steps of retrieval-only training), so seeds $0$--$2$ entered the abstention phase already retrieving
at RECUP $0.7734$, $0.7970$ and $0.7725$, while seeds $3$--$8$ started from scratch. The campaign's
pre-registration asserted that only the seed varied; the check performed at the time compared the
saved \emph{configs}, which do not record whether seeding occurred.

\begin{center}
\begin{tabular}{lccc}
\toprule
group & seeds & starting RECUP & outcome \\
\midrule
pre-trained & 0, 1, 2 & $\approx 0.78$ & \textbf{3 of 3 useful} \\
from scratch & 3--8 & $0$ & 2 useful, \textbf{4 degenerate} \\
\bottomrule
\end{tabular}
\end{center}

We state plainly what this costs and what it does not. \textbf{It costs the rate.} ``Four of nine
seeds'' is not a rate for a homogeneous population; the honest statement is $0$ of $3$ with
pre-training and $4$ of $6$ without. \textbf{It does not cost the phenomenon}, which is measured per
unit and does not depend on the grouping: the degenerate fixed point exists, four units are in it,
their heads sit at the prior, and the state is absorbing.

And it \emph{strengthens} the mechanism of \S\ref{sec:objetivo} rather than weakening it. If the
outcome is a race between retrieval learning and head settling, then a model that arrives already
retrieving has won the race before it starts---which is exactly the pattern observed. That two of the
six from-scratch units reach the useful point anyway shows pre-training is not necessary, only
strongly favourable.

\textbf{Where this does bite, and where it does not, is spelled out in the two affected sections}:
\S\ref{sec:exclusiva} rests on a paired control that shares both seed \emph{and} base and is
therefore unaffected, while the wider $31$-versus-$9$ tabulation is confounded and we do not rely on
it.

\subsection{Four of nine units end mute}
\label{sec:mudas}

\begin{table}[h]
\centering
\caption{The nine units at their own operating point. RECUP is retrieval ignoring the head; the last
three columns are the head's behaviour. Degenerate units in bold.}
\begin{tabular}{lrrrrrr}
\toprule
unit & step & RECUP & $\mathrm{AUC}(a)$ & $a > 0$ & median $a$ & range of $a$ \\
\midrule
b3\_s0 & 26000 & 0.9996 & 0.9997 & 0.407 & $-7.634$ & 29.17 \\
b3\_s1 & 26000 & 0.9992 & 0.9999 & 0.407 & $-7.563$ & 29.86 \\
b3\_s4 & 19000 & 0.7984 & 0.8084 & 0.413 & $-0.157$ & 14.17 \\
b3\_s2 & 26000 & 0.7946 & 0.8163 & 0.392 & $-0.233$ & 14.23 \\
b3\_s5 & 6000  & 0.7314 & 0.7550 & 0.529 & $+0.069$ & 5.66 \\
\midrule
\textbf{b3\_s6} & 26000 & \textbf{0.3964} & \textbf{0.5734} & \textbf{1.000} & $+1.006$ & \textbf{3.12} \\
\textbf{b3\_s3} & 26000 & \textbf{0.3820} & \textbf{0.5784} & \textbf{1.000} & $+1.045$ & \textbf{2.83} \\
\textbf{b3\_s7} & 13500 & \textbf{0.3432} & \textbf{0.5220} & \textbf{1.000} & $+1.251$ & \textbf{1.06} \\
\textbf{b3\_s8} & 8000  & \textbf{0.3040} & \textbf{0.5458} & \textbf{1.000} & $+1.507$ & \textbf{0.85} \\
\bottomrule
\end{tabular}
\label{tab:nueve}
\end{table}

\paragraph{What the standard metrics say about the bold rows.} All four abstain on every question.
Therefore \texttt{nose} $=1.0000$: they correctly decline every unanswerable question. And
\texttt{invento} $=0.0000$: they never name content absent from the archive, because they never name
anything. On the two numbers the abstention literature reports, these are the best units in the
table---strictly better than \texttt{b3\_s0}, which answers and therefore occasionally errs.

Their global accuracy is $0.4065$: the trivial floor, to four decimals, in all four.

\paragraph{A hypothesis we held and disproved.} \texttt{vigente} is composite, so
$\texttt{vigente} = 0.0000$ is compatible with \emph{perfect} retrieval hidden behind a stuck head.
We expected that. It is false: RECUP is $0.30$--$0.40$ in the degenerate units against $0.73$--$1.00$
in the rest. Retrieval is genuinely impaired.

But it is not zero, and that matters for the interpretation. At $\approx 0.35$ against a chance rate
below $0.01$, these models have learned a substantial part of the task and stopped short. The
degenerate state is not an untrained model; it is a partially trained model whose head correctly
reports that partial training as insufficient.

\subsection{The head has collapsed to the prior, and the constant is where theory puts it}
\label{sec:prior}

The four degenerate heads do not discriminate: $\mathrm{AUC}$ on their own target is
$0.522$--$0.578$, against $0.9997$ for \texttt{b3\_s0}. Nor are they saturated at $+\infty$: the
entire $1$--$99$ percentile range of their logits spans $0.85$--$3.24$, against $29.17$. They are
constants sitting just above the decision threshold.

Equation~\eqref{eq:prior} predicts \emph{which} constant, with no fitted parameter:

\begin{table}[h]
\centering
\caption{Closed-form prior against measured median logit. No parameters are fitted.}
\begin{tabular}{lrrrrr}
\toprule
unit & RECUP & $\mathbb{E}[b]$ & predicted $\mathrm{logit}$ & measured median $a$ & difference \\
\midrule
b3\_s8 & 0.3040 & 0.8196 & $1.513$ & $1.507$ & $\mathbf{0.006}$ \\
b3\_s7 & 0.3432 & 0.7963 & $1.363$ & $1.251$ & $0.112$ \\
b3\_s3 & 0.3820 & 0.7733 & $1.227$ & $1.045$ & $0.182$ \\
b3\_s6 & 0.3964 & 0.7647 & $1.179$ & $1.006$ & $0.173$ \\
\bottomrule
\end{tabular}
\label{tab:prior}
\end{table}

The ordering is reproduced in all four and the youngest unit falls within $0.006$ of the
parameter-free prediction. The older units sit slightly below their prior, which is what a head
beginning to discriminate would do; consistent with this, the range of $a$ grows monotonically with
step across the four ($0.85 \to 1.06 \to 2.83 \to 3.12$). \textbf{The head is not broken. It is
optimal for a generator that is wrong $80\%$ of the time.}

\subsection{The bifurcation is decided at step 2500}
\label{sec:bifurcacion}

Every unit with this target passes through an initial mute phase, which is expected: at
initialization the target is the constant~$1$. What differs is whether it ends.

The training loop already evaluates on $8$ batches of $64$, i.e.\ $512$ samples. At step $2500$:

\begin{table}[h]
\centering
\caption{Answers emitted out of $512$ at step $2500$, against the eventual outcome.}
\begin{tabular}{lrrl}
\toprule
unit & \texttt{abstencion} @2500 & answers emitted & outcome \\
\midrule
b3\_s2 & 0.9844 & 8 & useful \\
b3\_s0 & 0.9922 & 4 & useful \\
b3\_s5 & 0.9922 & 4 & useful \\
b3\_s1 & 0.9961 & 2 & useful \\
b3\_s4 & 0.9980 & \textbf{1} & useful \\
\midrule
b3\_s3 & 1.0000 & \textbf{0} & degenerate \\
b3\_s6 & 1.0000 & \textbf{0} & degenerate \\
b3\_s7 & 1.0000 & \textbf{0} & degenerate \\
b3\_s8 & 1.0000 & \textbf{0} & degenerate \\
\bottomrule
\end{tabular}
\label{tab:bifurcacion}
\end{table}

The separation is zero versus non-zero, and \texttt{b3\_s4} survives on a \textbf{single sample out
of $512$} and reaches $\texttt{vigente} = 0.68$. Breaking the silence once appears to be the event
that matters, which is what the race in \S\ref{sec:objetivo} predicts: one sample on which the
\texttt{argmax} is right is one sample whose target is $0$, and that is what a constant head has no
gradient to produce.

\paragraph{Validation on the repository.} Applied to \textbf{every} run in this repository with a
binary abstention head that reached $\geq 6000$ steps and has a step-$2500$ checkpoint---$40$ runs,
twelve families, two decision rules, both targets---the rule separates \textbf{40 of 40}: $36$
predicted useful and useful, $4$ predicted degenerate and degenerate.

\paragraph{The hardest available test, run after the rule was formulated.} The rule was derived
post-hoc. We then took the two \emph{oldest} degenerate units to full horizon under the original
flags, pre-registering the prediction that neither would leave total abstention. Neither did:
\texttt{b3\_s3} and \texttt{b3\_s6} reached step $26000$ with $\texttt{abstencion} \geq 0.999$ at
every logged milestone. This is not a prospective validation---the properly prospective test would
have been a screening design we did not run---but it is the strongest test the existing units
permitted, and the rule survived it.

\paragraph{A caveat on the margin.} One sample in $512$ is a margin of $0.2\%$. As an operational
criterion the abstention \emph{rate} is too fragile; the robust form of the same quantity is
$\min_i a_i$ over a large batch, which is continuous and independent of batch size.

\subsection{The mute phase belongs to the self-referential target}
\label{sec:exclusiva}

\textbf{This is the section the seeding confound touches most, so we state the limit first.} Every
run with a fixed target in this repository uses seeds $0$--$2$, and those are exactly the seeds for
which a pre-trained base exists. The wide tabulation below is therefore confounded with
initialization, and we report it as description rather than as evidence:

\begin{center}
\begin{tabular}{lrl}
\toprule
target & runs & \texttt{abstencion} at step 2500 \\
\midrule
fixed (``does the answer exist?'') & 31 & $0.02$--$0.32$ \\
\textbf{self-referential (``will I be wrong?'')} & 9 & $\mathbf{0.98}$--$\mathbf{1.00}$ \\
\bottomrule
\end{tabular}
\end{center}

\textbf{The claim rests instead on the paired control, which the confound does not reach.} The three
\texttt{p3} units share seed, architecture, budget \emph{and} pre-trained base with
\texttt{b3\_s0/s1/s2}, and differ only in the target:

\begin{center}
\begin{tabular}{lcc}
\toprule
seed & \texttt{abstencion} @2500, fixed target & \texttt{abstencion} @2500, self-referential \\
\midrule
0 & 0.086 & \textbf{0.992} \\
1 & 0.201 & \textbf{0.996} \\
2 & 0.063 & \textbf{0.984} \\
\bottomrule
\end{tabular}
\end{center}

Three of three, with the confounding variable held fixed by construction. The self-referential target
drives the model into the mute phase \emph{even when it arrives already retrieving} at RECUP
$\approx 0.78$; the fixed target never does. The \texttt{p3} units end at RECUP $0.80$--$0.97$ and
none falls in the degenerate band.

The mechanism is the one \S\ref{sec:objetivo} names: a fixed target is observable from step~$1$ and
discriminative immediately, so its head never passes through the constant phase in which the
degenerate solution is available.

\subsection{The degenerate point is absorbing, not slow}
\label{sec:absorbente}

Whether the state is a fixed point or a slow bottleneck changes what follows from it entirely, and
the two are distinguishable by measuring retrieval \emph{inside one unit}---which this line of work
had never done, because intermediate checkpoints are overwritten and the logged
\texttt{vigente} is $0.0000$ at every milestone regardless of retrieval.

We archived every partial checkpoint of \texttt{b3\_s3} from $22000$ to $26000$ and measured RECUP on
each with paired samples ($n=8000$, seed $54321$), under a pre-registration frozen beforehand that
asked for $\geq +0.0100$ and monotonicity in $\geq 3$ of $4$ intervals:

\begin{center}
\begin{tabular}{lrrrr}
\toprule
step & 22000 & 23000 & 24000 & 25000 \quad 26000 \\
\midrule
RECUP & 0.3841 & 0.3778 & 0.3844 & 0.3812 \quad 0.3820 \\
\bottomrule
\end{tabular}
\end{center}

The total is $\mathbf{-0.0021}$ over $4000$ steps. With $p_{\mathrm{flip}} = 0.1526$ measured, the
paired standard deviation is $0.0057$, so the total is $0.4\sigma$ and each interval is $0.1$--$1.2\sigma$:
\textbf{noise around a flat line}. Monotonicity was $2$ of $4$. The pre-registration's risk criterion
fired and its abandonment clause applied; a planned follow-up campaign of $70000$ steps was
cancelled.

\paragraph{What this does not separate, declared in advance.} The final $4000$ steps run with the
learning rate at the floor of its cosine schedule. ``The attractor is absorbing'' and ``nothing moves
at floor learning rate'' are both consistent with a flat curve, and the cancelled campaign was what
would have separated them. We therefore state the result as \emph{absorbing under this budget and
this schedule}, not as absorbing in general. This is the principal open question left by the paper.

\paragraph{One interval points the other way.} \texttt{b3\_s6}, which had only $1000$ steps
remaining, gives $+0.0080$ ($1.6\sigma$). The pre-registration had already ruled that a single
interval does not constitute a trend, and it does not, but it is the one measurement in this paper
pushing in the opposite direction and we report it as such.

\subsection{A correction to our own first reading}
\label{sec:correccion}

Our first analysis fitted a line through the four degenerate units ordered by step and extrapolated
that RECUP would reach $0.70$ at $\approx 84000$ steps---i.e.\ that the state was slow rather than
absorbing. Two checks were run before publishing that reading, and the second overturned it.

\paragraph{The check that passed.} The initial measurement used $n=2000$, whose $\sigma = 0.0135$ is
comparable to the between-unit spread, so we suspected the trend was measurement noise. Re-measuring
all four at $n=8000$ gives $r = 0.9864$ with a range of $12.5\sigma$. \textbf{The trend is real and
the suspicion was wrong.}

\paragraph{The check that failed.} The trend is real but not linear. Per-interval slopes decrease
monotonically---$+0.0071$, $+0.0048$, $+0.0014$ per $1000$ steps between successive units, and
$-0.0005$ within \texttt{b3\_s3}---and a saturating fit gives residual sum of squares $12.6\times$
smaller than the linear one. The extrapolation was an artifact of fitting a straight line to a
concave curve.

\paragraph{What we do \emph{not} conclude from the saturating fit.} It has three parameters for four
points, and its own prediction at step $26000$ ($0.4535$) misses the measured value ($0.3820$) by
$0.07$. It is not evidence. The claim in \S\ref{sec:absorbente} rests on the within-unit measurement,
which has no between-seed confound, and the saturating fit only explains why the earlier
extrapolation was wrong.

\section{What we do not claim}
\label{sec:noclaim}

\paragraph{This is not evidence that the self-referential target is useless.} Five of nine units
reach a useful state, two of them with $\mathrm{RECUP} \geq 0.999$ and $\mathrm{AUC} \geq 0.9999$.
The claim is that the objective admits a second stable solution and that a substantial minority of
seeds find it---not that the objective fails.

\paragraph{The degenerate units are not miscalibrated.} It would be convenient to call them broken.
They are not: their heads are optimal for their generators, to within $0.006$ of the parameter-free
prediction in the cleanest case. The failure is that being well calibrated about knowing nothing is
worth nothing, which is a statement about the objective and the metrics, not about calibration.

\paragraph{We do not claim the model knows when it does not know.} Everything here is supervised: the
head is trained against labels, and in the self-referential case against labels derived from
supervised comparison. The degenerate units in particular are a caution against that reading, since
they exhibit the outward signature of knowing when they do not know while knowing nothing at all.

\paragraph{The step-2500 rule is post-hoc.} It was formulated after observing outcomes, validated on
$40$ existing runs and tested on two units taken to full horizon under pre-registration. That is
strong retrospective evidence and one hard confirmation; it is not a prospective validation, and we
do not report it as one.

\section{Limitations}
\label{sec:limits}

\begin{itemize}
\item \textbf{The starting point was not controlled, and we found it late.} Seeds $0$--$2$ inherited
$12{,}000$ steps of retrieval pre-training and seeds $3$--$8$ did not, so the headline split is not a
rate over a homogeneous population (\S\ref{sec:mudas}). The per-unit phenomena are unaffected and the
paired control of \S\ref{sec:exclusiva} is clean, but any future campaign on this question must fix
the initialization across all units, and the harness must record whether seeding occurred---the saved
config does not, which is why the original check missed it.
\item \textbf{Scale.} $863{,}730$ parameters, a synthetic $242$-token language, one difficulty level,
$p_{\mathrm{nose}}$ fixed at $0.4$. Nothing here transfers automatically. We note only that the
degenerate fixed point is a property of the objective's algebra rather than of the model's size,
which makes the transfer question worth asking rather than answered.
\item \textbf{Budget and schedule are confounded with the absorbing claim}, as declared in
\S\ref{sec:absorbente}. This is the single largest gap in the paper.
\item \textbf{Nine units.} The $5$--$4$ split is a count, not a rate estimate; we report it as
observed and do not attach a confidence interval to a proportion from nine seeds.
\item \textbf{Two units are incomplete.} \texttt{b3\_s7} and \texttt{b3\_s8} were stopped at $13500$
and $8000$ steps after the campaign's blocking criterion became arithmetically unreachable. Their
degenerate status at those steps is measured, not extrapolated, but they did not run to horizon, and
the deviation is declared in the repository.
\item \textbf{One decision rule.} All units use the separate binary head. Whether the token or slot
interfaces admit the same fixed point under a self-referential target is untested.
\item \textbf{Directed literature search, not a systematic review.}
\end{itemize}

\section{Conclusion}

A hallucination detector wants to predict its own model's errors, so the natural target is
self-referential. That target has two stable fixed points, and the degenerate one is not a
pathology to be engineered away but a correct solution: a model that retrieves almost nothing and
reports, accurately, that it retrieves almost nothing. Four of the six units that began from scratch
found it; none of the three that began already retrieving did.

What makes this more than a curiosity is that the degenerate solution is \textbf{invisible to the
metrics the field reports}. It scores $\texttt{nose} = 1.0000$ and $\texttt{invento} = 0.0000$,
better than the units that work, while its global accuracy sits at the trivial floor to four decimals.
Any criterion built from abstention rate and fabrication rate is maximized by it. The remedy is not
subtle---report global accuracy against its trivial floor, and require it alongside any abstention
metric---but it has to be applied, and in this project it was applied only after a campaign had
already been designed around criteria that a mute model would have won.

The bifurcation is legible early and cheaply. Zero answers emitted at step $2500$ separated $40$ of
$40$ runs, and one unit survived on one sample out of $512$. If that rule generalizes---which this
paper does not establish---then the relevant question for a self-referential abstention objective is
not how to make a model abstain well. Abstaining well is trivial, and a model can do it by knowing
nothing. The question is how to leave silence without beginning to invent, and that is decided in the
first few thousand steps.

\section*{Data and code availability}

All pre-registrations with their SHA hashes, deviation files, analysis scripts, per-seed raw results
and machine-generated reports are archived in the project repository,
\url{https://github.com/SperanzaMax/delta-archivo}, whose commit history carries the timestamps that
establish the order of pre-registration and result.

\section*{Declarations}

\textbf{Funding.} None. \textbf{Competing interests.} None. \textbf{Ethics approval.} Not applicable;
no human or animal subjects. \textbf{Data availability.} See above. \textbf{Author contributions.}
Sole author.

\bibliographystyle{plainnat}
\begin{thebibliography}{9}

\bibitem[Arazo et~al.(2020)]{arazo2020pseudo}
Arazo, E., Ortego, D., Albert, P., O'Connor, N.~E., and McGuinness, K. (2020).
Pseudo-labeling and confirmation bias in deep semi-supervised learning.
\emph{IJCNN}.

\bibitem[Chen and He(2021)]{chen2021simsiam}
Chen, X. and He, K. (2021).
Exploring simple Siamese representation learning.
\emph{CVPR}.

\bibitem[Cohen et~al.(2024)]{cohen2024idk}
Cohen, R., Dobler, K., Biran, E., and de~Melo, G. (2024).
I don't know: Explicit modeling of uncertainty with an [IDK] token.
\emph{NeurIPS}.

\bibitem[El-Yaniv and Wiener(2010)]{elyaniv2010selective}
El-Yaniv, R. and Wiener, Y. (2010).
On the foundations of noise-free selective classification.
\emph{JMLR}, 11:1605--1641.

\bibitem[Geifman and El-Yaniv(2017)]{geifman2017selective}
Geifman, Y. and El-Yaniv, R. (2017).
Selective classification for deep neural networks.
\emph{NeurIPS}.

\bibitem[Geifman and El-Yaniv(2019)]{geifman2019selectivenet}
Geifman, Y. and El-Yaniv, R. (2019).
SelectiveNet: A deep neural network with an integrated reject option.
\emph{ICML}.

\bibitem[Grill et~al.(2020)]{grill2020byol}
Grill, J.-B. et~al. (2020).
Bootstrap your own latent: A new approach to self-supervised learning.
\emph{NeurIPS}.

\bibitem[Speranza(2026)]{speranza2026abstention}
Speranza, M. (2026).
Where abstention lives: A pre-registered four-way comparison of abstention interfaces in a small
language model with versioned memory.
\emph{Research Square}, DOI 10.21203/rs.3.rs-10839567/v1.

\end{thebibliography}

\end{document}
